% !TEX root = task-6-9-3.tex
% Editable task bank. Append new BankTask blocks at the end of the document.
\documentclass[12pt,a4paper]{article}
\usepackage{physics-textbook}
\usepackage{graphicx,needspace}
\usetikzlibrary{decorations.pathmorphing}
\newif\ifShowSolutions
\ifdefined\TutorEdition
  \ShowSolutionstrue
\else
  \ShowSolutionsfalse % Change to \ShowSolutionstrue to display solutions.
\fi
\newcounter{banktask}
\newsavebox{\BankTaskBox}
\newcommand{\BankTaskContent}[3]{%
  \noindent\textbf{\thebanktask.}\quad #2\par
  \ifShowSolutions
    \Needspace{50mm}% Keep the solution heading with its opening text or diagram.
    #3\par
    \ifdefined\TutorEdition
      \ifcsname TutorSource#1\endcsname
        \nopagebreak\smallskip{\small\csname TutorSource#1\endcsname\par}%
      \fi
    \fi
  \fi
}
\newcommand{\BankTask}[3]{%
  \refstepcounter{banktask}\par\medskip
  \sbox{\BankTaskBox}{\begin{minipage}{\linewidth}%
    \BankTaskContent{#1}{#2}{#3}\end{minipage}}%
  \ifdim\dimexpr\ht\BankTaskBox+\dp\BankTaskBox\relax>\dimexpr\textheight-20mm\relax
    \Needspace{35mm}\BankTaskContent{#1}{#2}{#3}%
  \else
    \Needspace{\dimexpr\ht\BankTaskBox+\dp\BankTaskBox\relax}%
    \noindent\usebox{\BankTaskBox}\par
  \fi
  \medskip
}
\begin{document}
\BookSubsection{6.9.3}{Закон радиоактивного распада}


% Source sdamgia-2109: https://phys-ege.sdamgia.ru/problem?id=2109
\BankTask{sdamgia-2109}{
Какая доля радиоактивных ядер распадается через интервал времени, равный половине периода полураспада? Ответ приведите в процентах и округлите до целых.
}{
\smallskip\noindent\textbf{Решение}\par
Согласно закону радиоактивного распада, по истечении времени $t= \frac{T}{2}$ от первоначального количества радиоактивных атомов $N_0$ останется примерно: 

\[N=N_0\cdot 2 ^{- t/T} =N_0\cdot 2 ^{- 1/2} = \frac{N_0}{\sqrt{2}} .\]

Тем самым доля ядер, не претерпевших радиоактивного распада, составит $\frac{N}{N_0} \cdot  100\%= \frac{1}{\sqrt{2}} \cdot  100\%.$ Доля же распавшихся ядер, очевидно, окажется равной $( 1 - \frac{1}{\sqrt{2}} ) \cdot  100\% \approx 29\%.$

Ответ: $29$.
}

% Source sdamgia-2124: https://phys-ege.sdamgia.ru/problem?id=2124
\BankTask{sdamgia-2124}{
Период полураспада ядер радиоактивного изотопа висмута $19\,\text{мин}$. Через какое время распадется $75\%$ ядер висмута в исследуемом образце? (Ответ дать в минутах.)
}{
\smallskip\noindent\textbf{Решение}\par
Согласно закону радиоактивного распада, по истечении времени $t$ от первоначального количества радиоактивных атомов $N_0$ останется примерно $N=N_0\cdot 2 ^{- t/T} .$ Таким образом, $75\%$ ядер висмута распадется (то есть $25\%$ останется) через время: 

\[t= - T \log_2  \frac{N}{N_0} =T \log_2  \frac{N_0}{N} =19\,\text{мин} \cdot  \log_2  \frac{1}{0{,}25} =38\,\text{мин}.\]

Ответ: $38$.
}

% Source sdamgia-2125: https://phys-ege.sdamgia.ru/problem?id=2125
\BankTask{sdamgia-2125}{
Период полураспада изотопа натрия ${}^{22}_{11}$Na равен $2{,}6$ года. Если изначально было $104\,\text{г}$ этого изотопа, то сколько примерно его будет через $5{,}2$ года? (Ответ дать в граммах.)
}{
\smallskip\noindent\textbf{Решение}\par
Согласно закону радиоактивного распада, по истечении времени $t$ от первоначальной массы радиоактивного изотопа $m_0$ останется примерно:

\[m=m_0\cdot 2^{-t/T}=104\,\text{г}\cdot 2^{-\frac{5{,}2\,\text{года}}{2{,}6\,\text{года}}}=26\,\text{г}.\]

Ответ: $26$.
}

% Source sdamgia-7192: https://phys-ege.sdamgia.ru/problem?id=7192
\BankTask{sdamgia-7192}{
Образец радиоактивного радия находится в закрытом сосуде. Ядра радия ${}^{224}_{88}\mathrm{Ra}$ испытывают $\alpha $-распад с периодом полураспада $3{,}6$ суток. Определите количество радия (в моль) в сосуде через $3{,}6$ суток, если в начальный момент времени образец содержал $1{,}8$ моль радия-$224$.
}{
\smallskip\noindent\textbf{Решение}\par
Период полураспада — это время, за которое распадается половина от всех имеющихся изначально атомов. 

Изначально имелось $1{,}8$ моль радия, значит, через период полураспада распадется $1{,}8$ : $2 = 0{,}9$ моль.

Ответ: $0{,}9$.
}

% Source sdamgia-7797: https://phys-ege.sdamgia.ru/problem?id=7797
\BankTask{sdamgia-7797}{
Изотоп технеция ${}^{95}_{43}\mathrm{Tc}$ испытывает позитронный $\beta $-распад с периодом полураспада $60$ суток, превращаясь в стабильный изотоп молибдена. В запаянную пробирку поместили $760\,\text{мг}$ указанного изотопа технеция. Сколько миллимолей технеция останется в пробирке через $180$ суток после начала опыта?
}{
\smallskip\noindent\textbf{Решение}\par
Найдем массу технеция в пробирке через $180$ суток: $m = m_0 \cdot  2 ^{- 180/60} =760\,\text{мг}/ 8 = 95\,\text{мг}.$ Следовательно, в пробирке останется $\nu = \frac{m}{\mu} = \frac{95\,\text{мг}}{95\,\text{г}/\text{моль}} = 1\,\text{ммоль}$ технеция.

Ответ: $1$.
}

% Source sdamgia-9747: https://phys-ege.sdamgia.ru/problem?id=9747
\BankTask{sdamgia-9747}{
После крупной радиационной аварии, произошедшей в $1986$ году на Чернобыльской атомной электростанции, некоторые участки местности оказались сильно загрязнены радиоактивным изотопом цезия-$137\,\text{с}$ периодом полураспада $30$ лет. На некоторых участках норма максимально допустимого содержания цезия-$137$ была превышена в $1000$ раз. Через сколько периодов полураспада после загрязнения такие участки местности вновь можно считать удовлетворяющими норме? Ответ округлите до целого числа.
}{
\smallskip\noindent\textbf{Решение}\par
В соответствии с законом радиоактивного распада получаем:

\[m=m_0 \cdot  2 ^{- t/T} \Longleftrightarrow  \frac{t}{T} = \log_2  \frac{m_0}{m} = \log_2  1000 \approx 10.\]

Загрязненные участки местности можно считать удовлетворяющими норме только спустя $10$ периодов полураспада.

Ответ: $10$.
}

% Source sdamgia-10192: https://phys-ege.sdamgia.ru/problem?id=10192
\BankTask{sdamgia-10192}{
Период полураспада элемента $1$ в три раза больше периода полураспада элемента $2$. За некоторое время число атомов элемента $1$ уменьшилось в $8$ раз. Во сколько раз за это же время уменьшилось число атомов элемента $2$?
}{
\smallskip\noindent\textbf{Решение}\par
Период полураспада — это время, в течение которого распадается половина наличного числа радиоактивных атомов.

Согласно закону радиоактивного распада, по истечении времени $t$ от первоначального количества атомов $N_0$ радиоактивного вещества с периодом полураспада T останется примерно $N=N_0\cdot 2 ^{- \frac{t}{T}}$ атомов. 

Поскольку число атомов элемента $1$ уменьшилось в $8$ раз, то прошло $3$ периода полураспада. За это же время число атомов элемента $2$ уменьшится в $2 ^{9} =512$ раз. 

Ответ: $512$.
}

% Source sdamgia-10325: https://phys-ege.sdamgia.ru/problem?id=10325
\BankTask{sdamgia-10325}{
Определите отношение числа распавшихся ядер некоторого радиоактивного изотопа к числу нераспавшихся ядер через время, равное семи периодам полураспада этого изотопа.
}{
\smallskip\noindent\textbf{Решение}\par
Согласно закону радиоактивного распада, по истечении времени $t$ от первоначального количества ядер $N_0$ радиоактивного вещества с периодом полураспада T останется примерно $N=N_0 \cdot  2 ^{- t/T}$ ядер. 

Тогда отношение распавшихся ядер к числу нераспавшихся равно:

\[\frac{N_0 - N}{N} = \frac{N_0}{N} - 1=2 ^{t} /T - 1=2 ^{7} - 1=127.\]

Ответ: $127$.
}

% Source sdamgia-11675: https://phys-ege.sdamgia.ru/problem?id=11675
\BankTask{sdamgia-11675}{
В пробирке в момент времени $t_{0}$ = $0$ находилось некоторое количество ядер радиоактивного изотопа. Через $t_{1}$ = $5\,\text{мин}$. в пробирке осталось $3416$ мкмоль нераспавшихся ядер, а через $t_{2}$ = $17\,\text{мин}$. — $427$ мкмоль нераспавшихся ядер. Чему равен период полураспада исходного изотопа? Ответ приведите в минутах.
}{
\smallskip\noindent\textbf{Решение}\par
Период полураспада — это время, в течение которого распадается половина начального числа радиоактивных атомов.

Согласно закону радиоактивного распада по истечении времени $t$ от первоначального количества атомов $N_0$ радиоактивного вещества с периодом полураспада T останется примерно $N=N_0\cdot 2 ^{- t/T}$ атомов. Аналогичное уравнение можно записать и для количества вещества: $\nu=\nu_0 2 ^{- t/T} .$ Запишем отношение количества вещества в пробирке через $5$ минут, к количеству вещества в пробирке через $17$ минут:

\[\frac{\nu_1}{\nu_2}=\frac{2^{-t_1/T}}{2^{-t_2/T}}=2^{-(t_1-t_2)/T}.\]

Откуда получаем: 

\[\frac{t_2 - t_1}{T} = \log_2  \frac{\nu_1}{\nu_2} \Longleftrightarrow  T = \frac{t_2 - t_1}{\log_2  \frac{\nu_1}{\nu_2}} .\]

Подставляя числовые выражения, находим: $T = \frac{17\,\text{мин} - 5\,\text{мин}}{\log_2  \frac{3416\,\text{мкмоль}}{427\,\text{мкмоль}}} = 4\,\text{мин}.$

Ответ: $4$.
}

% Source sdamgia-19673: https://phys-ege.sdamgia.ru/problem?id=19673
\BankTask{sdamgia-19673}{
Находящийся в пробирке радиоактивный таллий-$207$, начальное количество которого было равно $0{,}8$ моль, претерпевает радиоактивный распад, превращаясь в стабильный изотоп свинца. Какая масса свинца образуется в пробирке в течение трех периодов полураспада таллия, если молярная масса получающегося свинца равна $207\,\text{г}$/моль? Ответ выразите в граммах и округлите до целого числа.
}{
\smallskip\noindent\textbf{Решение}\par
Сначала найдем количество оставшегося по истечению указанного времени таллия:

\[\nu_{\text{ост}}\rm ( Tl ) = \frac{\nu_0 ( Tl )}{2 ^{\it t/T}} = \frac{\nu_0\rm ( Tl )}{8} = 0{,}1\,\text{моль}.\]

Тогда количество вещества образовавшегося свинца $\nu\rm ( Pb ) = \nu_0\rm ( Tl ) - \nu\rm ( Tl ) = 0{,}8 - 0{,}1 = 0{,}7\,\text{моль}.$ Отсюда масса свинца $m(\mathrm{Pb})=\nu(\mathrm{Pb})\cdot M(\mathrm{Pb})=0{,}7\cdot207=145\,\text{г}.$

Ответ: $145$.
}


% Source 4C2E4A: https://ege.fipi.ru/bank/questions.php?search=1&proj=BA1F39653304A5B041B656915DC36B38&qid=4C2E4A
\BankTask{4C2E4A}{
Период полураспада \(T\) изотопа селена \({}_{34}^{81}\mathrm{Se}\) равен \(18\,\text{мин}\). Какая масса (в \(\text{мг}\)) этого изотопа осталась в образце, содержавшем первоначально \(120\,\text{мг}\) \({}_{34}^{81}\mathrm{Se}\), через \(54\,\text{мин}\)?
}{
\smallskip\noindent\textbf{Решение от GPT}\par
За \(54\,\text{мин}\) проходит \(n=t/T=54/18=3\) периода полураспада. Оставшаяся масса \(m=m_0\cdot2^{-n}=120/8=15\,\text{мг}\). Ответ: \(15\,\text{мг}\).
}

% Source 453BF5: https://ege.fipi.ru/bank/questions.php?search=1&proj=BA1F39653304A5B041B656915DC36B38&qid=453BF5
\BankTask{453BF5}{
Дан график зависимости числа нераспавшихся ядер некоторого изотопа от времени. Каков период полураспада этого изотопа?
\begin{center}
\begin{tikzpicture}[book diagram,x=0.15cm,y=0.75cm]
\draw[thin] (0,0) grid[xstep=5,ystep=1] (35,7);
\draw[-{book tip}] (0,0) -- (37,0);
\draw[-{book tip}] (0,0) -- (0,7.85);
\foreach \x in {5,10,15,20,25,30,35}{\node[below] at (\x,0) {\(\x\)};}
\foreach \y in {1,2,3,4,5,6,7}{\node[left] at (0,\y) {\(\y\)};}
\node[below left] at (0,0) {\(0\)};
\node[anchor=east] at (-0.3,7.75) {\(N,\,10^{20}\)};
\node[anchor=west] at (37,0) {\(t,\,\text{мин.}\)};
\draw[thick,smooth,domain=0:30,samples=121] plot (\x,{6*pow(2,-\x/5)});
\end{tikzpicture}
\end{center}
}{
\smallskip\noindent\textbf{Решение от GPT}\par
По графику начальное число ядер равно \(N_0=6\cdot10^{20}\). Число ядер уменьшается вдвое, до \(3\cdot10^{20}\), за \(5\,\text{мин}\), поэтому период полураспада \(T=5\,\text{мин}\). Ответ: \(5\,\text{мин}\).
}

% Source 8F0C7B: https://ege.fipi.ru/bank/questions.php?search=1&proj=BA1F39653304A5B041B656915DC36B38&qid=8F0C7B
\BankTask{8F0C7B}{
Закон радиоактивного распада ядер некоторого изотопа имеет вид: \(N=N_0\cdot2^{-\lambda t}\), где \(\lambda=0{,}04\ \text{с}^{-1}\). Определите период полураспада этих ядер (в \(\text{с}\)).
}{
\smallskip\noindent\textbf{Решение от GPT}\par
За период полураспада \(N/N_0=1/2\), поэтому \(2^{-\lambda T}=2^{-1}\) и \(\lambda T=1\). Следовательно, \(T=1/\lambda=1/0{,}04=25\,\text{с}\). Ответ: \(25\,\text{с}\).
}

% Source AB2CB2: https://ege.fipi.ru/bank/questions.php?search=1&proj=BA1F39653304A5B041B656915DC36B38&qid=AB2CB2
\BankTask{AB2CB2}{
Образец радиоактивного радия \({}_{88}^{224}\mathrm{Ra}\) находится в закрытом сосуде, из которого откачан воздух. Ядра радия испытывают \(\alpha\)-распад с периодом полураспада \(3{,}6\,\text{суток}\). Определите число моль гелия в сосуде через \(10{,}8\,\text{суток}\), если образец в момент его помещения в сосуд имел в своём составе \(6\cdot10^{23}\) атомов радия-224, а атомов гелия в сосуде не было.
}{
\smallskip\noindent\textbf{Решение от GPT}\par
За время \(t=10{,}8\,\text{суток}\) проходит \(t/T=10{,}8/3{,}6=3\) периода полураспада, поэтому распадается \(N_0(1-2^{-3})=\frac78N_0\) ядер радия. При каждом \(\alpha\)-распаде образуется одно ядро гелия, поэтому \(\nu_{\mathrm{He}}=\frac{7N_0}{8N_A}=\frac{7\cdot6\cdot10^{23}}{8\cdot6\cdot10^{23}}=0{,}875\ \text{моль}\). Ответ: \(0{,}875\,\text{моль}\).
}

% Source BCED39: https://ege.fipi.ru/bank/questions.php?search=1&proj=BA1F39653304A5B041B656915DC36B38&qid=BCED39
\BankTask{BCED39}{
Чему равен период (в годах) полураспада ядер изотопа \({}_{11}^{22}\mathrm{Na}\), если за \(9\,\text{лет}\) исходно большое число нераспавшихся ядер атомов этого изотопа уменьшилось в \(8\) раз?
}{
\smallskip\noindent\textbf{Решение от GPT}\par
По закону радиоактивного распада \(N/N_0=2^{-t/T}=1/8=2^{-3}\), следовательно, \(t/T=3\). Период полураспада \(T=9/3=3\,\text{года}\). Ответ: \(3\,\text{года}\).
}

% Source E423E8: https://ege.fipi.ru/bank/questions.php?search=1&proj=BA1F39653304A5B041B656915DC36B38&qid=E423E8
\BankTask{E423E8}{
Через сколько лет из $2\cdot10^{10}$ ядер радиоактивного изотопа цезия ${}_{55}^{137}\mathrm{Cs}$, имеющего период полураспада $T=26\,\text{лет}$, нераспавшимися останутся $2{,}5\cdot10^9$ ядер изотопа?
}{
\smallskip\noindent\textbf{Решение от GPT}\par
По закону радиоактивного распада $N=N_0\cdot2^{-t/T}$. Отношение $N_0/N=(2\cdot10^{10})/(2{,}5\cdot10^9)=8=2^3$, поэтому прошло три периода полураспада. Следовательно, $t=3T=78\,\text{лет}$.
}

% Source DCA275: https://ege.fipi.ru/bank/questions.php?search=1&proj=BA1F39653304A5B041B656915DC36B38&qid=DCA275
\BankTask{DCA275}{
Какая доля (в процентах) от исходного большого количества радиоактивных ядер распадётся за интервал времени, равный двум периодам полураспада?
}{
\smallskip\noindent\textbf{Решение от GPT}\par
За два периода полураспада останется доля ядер $N/N_0=2^{-2}=1/4$. Доля распавшихся ядер равна $1-1/4=3/4=75\,\%$. Ответ: $75\,\%$.
}

% Source 38C6C8: https://ege.fipi.ru/bank/questions.php?search=1&proj=BA1F39653304A5B041B656915DC36B38&qid=38C6C8
\BankTask{38C6C8}{
Образец радиоактивного полония ${}_{84}^{218}\mathrm{Po}$ находится в закрытом сосуде, из которого откачан воздух. Ядра полония испытывают $\alpha$-распад с периодом полураспада $3\,\text{мин}$. Определите число моль полония-218 в сосуде через $9\,\text{мин.}$, если образец в момент его помещения в сосуд имел в своём составе $2{,}4\times10^{23}$ атомов полония-218.
}{
\smallskip\noindent\textbf{Решение от GPT}\par
За время $9\,\text{мин}$ проходит $t/T=9/3=3$ периода полураспада. Число оставшихся атомов полония-218 равно $N=N_0\cdot2^{-3}=3\times10^{22}$. Количество вещества $\nu=N/N_A=\dfrac{3\times10^{22}}{6\times10^{23}}\,\text{моль}=0{,}05\,\text{моль}$. Ответ: $0{,}05\,\text{моль}$.
}

\end{document}
